\documentclass[12pt, a4paper]{report}

% ── Encodage & langue ──────────────────────────────────────────────────────────
\usepackage[utf8]{inputenc}
\usepackage[T1]{fontenc}
\usepackage[francais]{babel}

% ── Mise en page ───────────────────────────────────────────────────────────────
\usepackage[top=2.5cm, bottom=2.5cm, left=3cm, right=2.5cm]{geometry}
\usepackage{setspace}
\onehalfspacing
\setlength{\parindent}{0pt}
\setlength{\parskip}{6pt}

% ── Typographie ────────────────────────────────────────────────────────────────
\usepackage{ae}

% ── Mathématiques ──────────────────────────────────────────────────────────────
\usepackage{amsmath, amssymb, amsthm}
\usepackage{mathtools}

% ── Tableaux ───────────────────────────────────────────────────────────────────
\usepackage{booktabs}
\usepackage{array}
\usepackage{multirow}
\usepackage{xcolor}
\usepackage{colortbl}
\usepackage{tabularx}
\usepackage{longtable}

% ── Graphiques & figures ───────────────────────────────────────────────────────
\usepackage{graphicx}
\usepackage{tikz}
\usepackage{pgfplots}
\pgfplotsset{compat=1.18}
\usetikzlibrary{shapes.geometric, arrows.meta, positioning, calc, fit, backgrounds}

% ── Listings (code) ────────────────────────────────────────────────────────────
\usepackage{listings}
\usepackage{lstautogobble}

% ── En-têtes / pieds de page ───────────────────────────────────────────────────
\usepackage{fancyhdr}
\pagestyle{fancy}
\fancyhf{}
\fancyhead[L]{\small\textit{Métriques YOLOv8}}
\fancyhead[R]{\small\textit{\leftmark}}
\fancyfoot[C]{\thepage}
\renewcommand{\headrulewidth}{0.4pt}

% ── Hyperliens ─────────────────────────────────────────────────────────────────
\usepackage[hidelinks, colorlinks=false]{hyperref}
\usepackage{cleveref}

% ── Boîtes colorées ────────────────────────────────────────────────────────────
\usepackage[most]{tcolorbox}
\tcbuselibrary{breakable, skins}

% ── Couleurs personnalisées ────────────────────────────────────────────────────
\definecolor{yoloblue}{RGB}{0, 82, 165}
\definecolor{yoloorange}{RGB}{230, 100, 20}
\definecolor{yologreen}{RGB}{34, 139, 34}
\definecolor{yolored}{RGB}{180, 30, 30}
\definecolor{lightgray}{RGB}{245, 245, 245}
\definecolor{lightblue}{RGB}{230, 241, 255}
\definecolor{lightorange}{RGB}{255, 243, 224}
\definecolor{lightgreen}{RGB}{230, 255, 230}
\definecolor{lightred}{RGB}{255, 230, 230}

% ── Environnements tcolorbox ───────────────────────────────────────────────────
\newtcolorbox{definitionbox}[1]{
  colback=lightblue, colframe=yoloblue, fonttitle=\bfseries,
  title={#1}, breakable, arc=4pt
}
\newtcolorbox{formulebox}[1]{
  colback=lightorange, colframe=yoloorange, fonttitle=\bfseries,
  title={#1}, breakable, arc=4pt
}
\newtcolorbox{examplebox}[1]{
  colback=lightgreen, colframe=yologreen, fonttitle=\bfseries,
  title={#1}, breakable, arc=4pt
}
\newtcolorbox{warningbox}[1]{
  colback=lightred, colframe=yolored, fonttitle=\bfseries,
  title={#1}, breakable, arc=4pt
}
\newtcolorbox{notebox}{
  colback=lightgray, colframe=gray!60, fonttitle=\bfseries,
  title={Note}, breakable, arc=4pt
}

% ── Style listings Python ──────────────────────────────────────────────────────
\lstdefinestyle{python}{
  language=Python,
  basicstyle=\ttfamily\small,
  keywordstyle=\color{yoloblue}\bfseries,
  commentstyle=\color{gray}\itshape,
  stringstyle=\color{yologreen},
  numberstyle=\tiny\color{gray},
  numbers=left,
  stepnumber=1,
  showstringspaces=false,
  breaklines=true,
  frame=single,
  rulecolor=\color{gray!40},
  backgroundcolor=\color{lightgray},
  autogobble=true
}

% ── Commandes utilitaires ──────────────────────────────────────────────────────
\newcommand{\mAP}{\text{mAP}}
\newcommand{\IoU}{\text{IoU}}
\newcommand{\tp}{\text{TP}}
\newcommand{\fp}{\text{FP}}
\newcommand{\fn}{\text{FN}}
\newcommand{\tn}{\text{TN}}

% ──────────────────────────────────────────────────────────────────────────────
%  DOCUMENT
% ──────────────────────────────────────────────────────────────────────────────
\begin{document}

% ════════════════════════════════════════════════════════════════════════════════
%  PAGE DE TITRE
% ════════════════════════════════════════════════════════════════════════════════
\begin{titlepage}
  \centering
  \vspace*{1cm}

  {\color{yoloblue}\rule{\linewidth}{2pt}}
  \vspace{0.5cm}

  {\Huge\bfseries\color{yoloblue}
    Métriques de Précision et d'Erreur\\[0.4cm]
    du Modèle \textcolor{yoloorange}{YOLOv8}
  }

  \vspace{0.4cm}
  {\color{yoloblue}\rule{\linewidth}{2pt}}
  \vspace{1.5cm}

  {\Large\textit{Détection et Reconnaissance d'Objets}}

  \vspace{1cm}

  \begin{tcolorbox}[colback=lightblue, colframe=yoloblue, width=0.85\linewidth,
    arc=6pt, boxrule=1pt]
    \centering
    \textbf{Référence exhaustive} — Définitions, formules, calculs détaillés\\
    et exemples numériques sur des données réelles
  \end{tcolorbox}

  \vspace{2cm}

  % Schéma pipeline simplifié
  \begin{tikzpicture}[node distance=1.4cm, >=Stealth,
    box/.style={rectangle, rounded corners=4pt, draw=yoloblue,
                fill=lightblue, minimum width=2.4cm, minimum height=0.9cm,
                font=\small\bfseries, text=yoloblue}]
    \node[box] (img)   {Image};
    \node[box, right=of img]  (yolo)  {YOLOv8};
    \node[box, right=of yolo] (det)   {Détections};
    \node[box, right=of det]  (eval)  {Évaluation};
    \node[box, right=of eval] (met)   {Métriques};
    \draw[->] (img)  -- (yolo);
    \draw[->] (yolo) -- (det);
    \draw[->] (det)  -- (eval);
    \draw[->] (eval) -- (met);
  \end{tikzpicture}

  \vfill

  {\large\today}
\end{titlepage}

% ════════════════════════════════════════════════════════════════════════════════
%  TABLE DES MATIÈRES
% ════════════════════════════════════════════════════════════════════════════════
\tableofcontents
\newpage

% ════════════════════════════════════════════════════════════════════════════════
%  INTRODUCTION
% ════════════════════════════════════════════════════════════════════════════════
\chapter{Introduction}

\section{Contexte et objectifs}

YOLOv8 (\textit{You Only Look Once version 8}), développé par Ultralytics, est
l'une des architectures les plus performantes pour la \textbf{détection et la
reconnaissance d'objets} en temps réel. Évaluer correctement un tel modèle
requiert de maîtriser un ensemble de métriques couvrant quatre dimensions
complémentaires :

\begin{itemize}
  \item \textbf{Précision de la classification} : le modèle attribue-t-il la
        bonne étiquette ?
  \item \textbf{Qualité de la localisation} : la boîte englobante est-elle bien
        positionnée ?
  \item \textbf{Erreurs d'entraînement} : les fonctions de perte convergent-elles
        correctement ?
  \item \textbf{Performance opérationnelle} : le modèle peut-il tourner en temps
        réel ?
\end{itemize}

\section{Notions fondamentales : TP, FP, FN, TN}

\begin{definitionbox}{Définitions de base}
  Soit une image contenant des objets réels (Ground Truth, GT).
  Le modèle produit des \textbf{prédictions}. On compare chaque prédiction au GT
  via un seuil d'IoU (typiquement 0{,}50) :
  \begin{itemize}
    \item \textbf{TP (True Positive)} : la prédiction correspond à un objet réel
          et la classe est correcte, \textit{et} $\IoU \geq \text{seuil}$.
    \item \textbf{FP (False Positive)} : la prédiction ne correspond à aucun objet
          réel, ou la classe est incorrecte.
    \item \textbf{FN (False Negative)} : un objet réel n'a pas été détecté.
    \item \textbf{TN (True Negative)} : zone sans objet, non détectée (rarement
          utilisé en détection).
  \end{itemize}
\end{definitionbox}

\begin{examplebox}{Exemple illustratif — 5 images, classe ``voiture''}
  \begin{center}
  \begin{tabular}{lccccc}
    \toprule
    \textbf{Image} & \textbf{GT réels} & \textbf{Détections} & \textbf{TP} & \textbf{FP} & \textbf{FN} \\
    \midrule
    Img-1 & 2 & 2 & 2 & 0 & 0 \\
    Img-2 & 3 & 4 & 3 & 1 & 0 \\
    Img-3 & 1 & 0 & 0 & 0 & 1 \\
    Img-4 & 2 & 2 & 1 & 1 & 1 \\
    Img-5 & 0 & 1 & 0 & 1 & 0 \\
    \midrule
    \textbf{Total} & \textbf{8} & \textbf{9} & \textbf{6} & \textbf{3} & \textbf{2} \\
    \bottomrule
  \end{tabular}
  \end{center}
\end{examplebox}

Les données de cet exemple ($\tp=6$, $\fp=3$, $\fn=2$) seront réutilisées dans
tout le document.

% ════════════════════════════════════════════════════════════════════════════════
%  CHAPITRE 1 — MÉTRIQUES DE PRÉCISION
% ════════════════════════════════════════════════════════════════════════════════
\chapter{Métriques de Précision et de Rappel}

% ─────────────────────────────────────────────────────────────────────────────
\section{Precision (P)}

\begin{definitionbox}{Définition}
  La \textbf{Precision} mesure la proportion de détections correctes parmi
  \emph{toutes} les détections émises par le modèle. Elle répond à la question :
  \og Quand le modèle dit qu'il a trouvé un objet, a-t-il raison ? \fg
\end{definitionbox}

\begin{formulebox}{Formule}
  \[
    P = \frac{\tp}{\tp + \fp}
  \]
  \begin{itemize}
    \item $\tp$ : nombre de vraies détections correctes
    \item $\fp$ : nombre de fausses alarmes (détections erronées)
  \end{itemize}
\end{formulebox}

\textbf{Interprétation :}
\begin{itemize}
  \item $P = 1{,}0$ : toutes les détections sont correctes, aucune fausse alarme.
  \item $P = 0{,}5$ : une détection sur deux est fausse.
  \item Une précision faible est dangereuse dans des systèmes de surveillance
        (trop d'alertes inutiles).
\end{itemize}

\begin{examplebox}{Calcul numérique}
  En reprenant les données du tableau ($\tp=6$, $\fp=3$) :
  \[
    P = \frac{6}{6 + 3} = \frac{6}{9} \approx \mathbf{0{,}667}
  \]
  \textit{Interprétation :} $66{,}7\,\%$ des détections émises par le modèle
  correspondent effectivement à des voitures réelles.
\end{examplebox}

% ─────────────────────────────────────────────────────────────────────────────
\section{Recall (R) — Rappel}

\begin{definitionbox}{Définition}
  Le \textbf{Recall} mesure la proportion d'objets réels effectivement détectés
  parmi tous ceux présents dans les images. Il répond à la question :
  \og Le modèle manque-t-il des objets ? \fg
\end{definitionbox}

\begin{formulebox}{Formule}
  \[
    R = \frac{\tp}{\tp + \fn}
  \]
  \begin{itemize}
    \item $\fn$ : objets réels non détectés
  \end{itemize}
\end{formulebox}

\textbf{Interprétation :}
\begin{itemize}
  \item $R = 1{,}0$ : aucun objet n'a été manqué.
  \item $R = 0{,}5$ : la moitié des objets n'est pas détectée.
  \item Un rappel faible est critique dans des applications médicales ou de
        sécurité (un objet manqué peut avoir de graves conséquences).
\end{itemize}

\begin{examplebox}{Calcul numérique}
  Avec $\tp=6$ et $\fn=2$ :
  \[
    R = \frac{6}{6 + 2} = \frac{6}{8} = \mathbf{0{,}75}
  \]
  \textit{Interprétation :} le modèle détecte $75\,\%$ des voitures réellement
  présentes dans les images.
\end{examplebox}

\begin{warningbox}{Compromis Precision / Recall}
  Il existe un \textbf{trade-off fondamental} : augmenter le seuil de confiance
  améliore la Precision (moins de fausses alarmes) mais réduit le Recall (plus
  d'objets manqués), et vice-versa. Le F1-Score et la courbe PR permettent
  d'analyser ce compromis.
\end{warningbox}

% ─────────────────────────────────────────────────────────────────────────────
\section{F1-Score}

\begin{definitionbox}{Définition}
  Le \textbf{F1-Score} est la moyenne harmonique de la Precision et du Recall.
  Il fournit un indicateur unique équilibrant les deux métriques précédentes.
\end{definitionbox}

\begin{formulebox}{Formule}
  \[
    F1 = 2 \times \frac{P \times R}{P + R}
  \]
  \textit{Équivalent développé :}
  \[
    F1 = \frac{2\,\tp}{2\,\tp + \fp + \fn}
  \]
\end{formulebox}

\textbf{Pourquoi la moyenne harmonique ?}

La moyenne arithmétique $(P+R)/2$ peut donner un score élevé même si l'une des
deux métriques est très faible. Par exemple, si $P=1{,}0$ et $R=0{,}1$ :
\[
  \text{Moy. arith.} = \frac{1{,}0 + 0{,}1}{2} = 0{,}55 \quad \text{vs} \quad
  F1 = \frac{2 \times 1{,}0 \times 0{,}1}{1{,}0 + 0{,}1} \approx 0{,}18
\]
La moyenne harmonique pénalise fortement les déséquilibres.

\begin{examplebox}{Calcul numérique}
  Avec $P = 0{,}667$ et $R = 0{,}75$ :
  \[
    F1 = 2 \times \frac{0{,}667 \times 0{,}75}{0{,}667 + 0{,}75}
       = 2 \times \frac{0{,}500}{1{,}417} \approx \mathbf{0{,}706}
  \]
  \textit{Vérification directe :}
  $F1 = \frac{2 \times 6}{2 \times 6 + 3 + 2} = \frac{12}{17} \approx 0{,}706$ \checkmark
\end{examplebox}

\textbf{Utilisation en YOLOv8 :} la courbe F1-Confidence permet de sélectionner
le seuil de confiance optimal (celui qui maximise le F1-Score sur les données de
validation).

% ─────────────────────────────────────────────────────────────────────────────
\section{AP — Average Precision}

\begin{definitionbox}{Définition}
  L'\textbf{Average Precision} est l'aire sous la courbe Precision-Recall (PR).
  Elle synthétise la performance du modèle sur une \emph{classe donnée} sur
  \emph{tous les seuils de confiance possibles}.
\end{definitionbox}

\begin{formulebox}{Formule (interpolation 11 points)}
  \[
    AP = \frac{1}{11} \sum_{r \in \{0,\,0.1,\,\ldots,\,1.0\}} P_{\text{interp}}(r)
  \]
  où $P_{\text{interp}}(r) = \max_{\tilde{r} \geq r} P(\tilde{r})$

  \textit{Ou, en intégrale continue (COCO) :}
  \[
    AP = \int_0^1 P(R)\, dR
  \]
\end{formulebox}

\subsection{Construction de la courbe PR}

\begin{enumerate}
  \item Trier toutes les détections par score de confiance décroissant.
  \item Pour chaque seuil de confiance $c_i$, calculer $P(c_i)$ et $R(c_i)$.
  \item Tracer la courbe $P = f(R)$.
  \item Calculer l'aire sous cette courbe.
\end{enumerate}

\begin{examplebox}{Calcul détaillé de l'AP}

  Supposons 10 détections triées par confiance décroissante pour la classe
  ``voiture'' (GT total = 5 objets). ``1'' = TP, ``0'' = FP :

  \vspace{0.3cm}
  \begin{center}
  \begin{tabular}{ccccc}
    \toprule
    \textbf{Rang} & \textbf{Conf.} & \textbf{TP/FP} & \textbf{Precision} & \textbf{Recall} \\
    \midrule
    1 & 0.95 & 1 & 1/1 = 1.000 & 1/5 = 0.200 \\
    2 & 0.90 & 1 & 2/2 = 1.000 & 2/5 = 0.400 \\
    3 & 0.85 & 0 & 2/3 = 0.667 & 2/5 = 0.400 \\
    4 & 0.80 & 1 & 3/4 = 0.750 & 3/5 = 0.600 \\
    5 & 0.75 & 0 & 3/5 = 0.600 & 3/5 = 0.600 \\
    6 & 0.70 & 1 & 4/6 = 0.667 & 4/5 = 0.800 \\
    7 & 0.65 & 0 & 4/7 = 0.571 & 4/5 = 0.800 \\
    8 & 0.60 & 1 & 5/8 = 0.625 & 5/5 = 1.000 \\
    9 & 0.50 & 0 & 5/9 = 0.556 & 5/5 = 1.000 \\
    10 & 0.40 & 0 & 5/10 = 0.500 & 5/5 = 1.000 \\
    \bottomrule
  \end{tabular}
  \end{center}

  Après interpolation (valeur max de $P$ pour $R \geq r$) aux 11 points :
  \[
    AP = \frac{1}{11}(1{,}0 + 1{,}0 + 1{,}0 + 0{,}75 + 0{,}75 + 0{,}667 +
         0{,}625 + 0{,}625 + 0 + 0 + 0)
       \approx \mathbf{0{,}584}
  \]
\end{examplebox}

% ─────────────────────────────────────────────────────────────────────────────
\section{mAP50 — mean Average Precision @ IoU=0.50}

\begin{definitionbox}{Définition}
  Le \textbf{mAP50} est la moyenne des AP de toutes les classes, calculée avec
  un seuil IoU fixé à $0{,}50$. C'est la \textbf{métrique principale} affichée
  par YOLOv8 lors de l'évaluation.
\end{definitionbox}

\begin{formulebox}{Formule}
  \[
    \mAP_{50} = \frac{1}{N} \sum_{i=1}^{N} AP_i^{(\IoU=0.50)}
  \]
  où $N$ est le nombre de classes.
\end{formulebox}

\begin{examplebox}{Calcul numérique — 4 classes}
  \begin{center}
  \begin{tabular}{lcc}
    \toprule
    \textbf{Classe} & \textbf{AP\textsubscript{50}} & \textbf{Observations} \\
    \midrule
    Voiture    & 0.892 & Classe bien représentée \\
    Piéton     & 0.741 & Petits objets, plus difficile \\
    Vélo       & 0.654 & Souvent occulté \\
    Camion     & 0.813 & Grands objets, bien détectés \\
    \midrule
    \textbf{mAP50} & $\mathbf{0.775}$ & $(0.892+0.741+0.654+0.813)/4$ \\
    \bottomrule
  \end{tabular}
  \end{center}
  \textit{Interprétation :} le modèle obtient en moyenne $77{,}5\,\%$ de précision
  sur les 4 classes avec un seuil de recoupement de 50\,\%.
\end{examplebox}

% ─────────────────────────────────────────────────────────────────────────────
\section{mAP50-95 — mean Average Precision @ IoU=0.50:0.95}

\begin{definitionbox}{Définition}
  Le \textbf{mAP50-95} est la moyenne des mAP calculés sur 10 seuils IoU
  distincts : $\{0.50,\, 0.55,\, 0.60,\, \ldots,\, 0.95\}$. C'est la métrique
  standard du challenge \textbf{COCO}, la plus exigeante.
\end{definitionbox}

\begin{formulebox}{Formule}
  \[
    \mAP_{50:95} = \frac{1}{10} \sum_{t \in \{0.50,\,0.55,\,\ldots,\,0.95\}}
    \mAP_t
  \]
\end{formulebox}

\begin{examplebox}{Calcul numérique}
  \begin{center}
  \begin{tabular}{cc}
    \toprule
    \textbf{Seuil IoU} & \textbf{mAP\textsubscript{t}} \\
    \midrule
    0.50 & 0.775 \\
    0.55 & 0.742 \\
    0.60 & 0.708 \\
    0.65 & 0.671 \\
    0.70 & 0.623 \\
    0.75 & 0.561 \\
    0.80 & 0.489 \\
    0.85 & 0.401 \\
    0.90 & 0.283 \\
    0.95 & 0.112 \\
    \midrule
    \textbf{mAP50-95} & $\mathbf{0.537}$ \\
    \bottomrule
  \end{tabular}
  \end{center}
  La chute des valeurs montre que le modèle se dégrade significativement quand
  on exige une localisation plus précise (IoU $\geq 0{,}75$).
\end{examplebox}

\begin{warningbox}{mAP50 vs mAP50-95}
  Un modèle peut afficher un excellent mAP50 ($\sim\!0{,}90$) mais un mAP50-95
  médiocre ($\sim\!0{,}45$), indiquant que les boîtes englobantes sont imprécises.
  Pour des applications critiques (robotique, véhicules autonomes), le mAP50-95
  est la métrique pertinente.
\end{warningbox}

% ════════════════════════════════════════════════════════════════════════════════
%  CHAPITRE 2 — MÉTRIQUES DE LOCALISATION
% ════════════════════════════════════════════════════════════════════════════════
\chapter{Métriques de Localisation — Bounding Box}

% ─────────────────────────────────────────────────────────────────────────────
\section{IoU — Intersection over Union}

\begin{definitionbox}{Définition}
  L'\textbf{IoU} mesure le taux de recoupement entre la boîte prédite
  $B_{\text{pred}}$ et la boîte réelle $B_{\text{gt}}$. C'est la base de toutes
  les métriques de localisation.
\end{definitionbox}

\begin{formulebox}{Formule}
  \[
    \IoU = \frac{\left|B_{\text{pred}} \cap B_{\text{gt}}\right|}
                {\left|B_{\text{pred}} \cup B_{\text{gt}}\right|}
  \]
  \textit{Développé :}
  \[
    \IoU = \frac{\text{Aire d'intersection}}
                {\text{Aire}(B_{\text{pred}}) + \text{Aire}(B_{\text{gt}}) - \text{Aire d'intersection}}
  \]
\end{formulebox}

\textbf{Calcul de l'intersection :}

Soient $B_{\text{pred}} = (x_1^p, y_1^p, x_2^p, y_2^p)$ et
$B_{\text{gt}} = (x_1^g, y_1^g, x_2^g, y_2^g)$ en coordonnées
coin-supérieur-gauche / coin-inférieur-droit.

\[
  x_{\text{inter},1} = \max(x_1^p, x_1^g), \quad
  x_{\text{inter},2} = \min(x_2^p, x_2^g)
\]
\[
  y_{\text{inter},1} = \max(y_1^p, y_1^g), \quad
  y_{\text{inter},2} = \min(y_2^p, y_2^g)
\]
\[
  \text{Aire}_{\text{inter}} =
    \max(0,\; x_{\text{inter},2} - x_{\text{inter},1}) \times
    \max(0,\; y_{\text{inter},2} - y_{\text{inter},1})
\]

\begin{examplebox}{Calcul numérique}
  \begin{itemize}
    \item $B_{\text{pred}} = (100,\,150,\,250,\,350)$ — Aire $= 150 \times 200 = 30\,000$
    \item $B_{\text{gt}} = (120,\,160,\,270,\,370)$ — Aire $= 150 \times 210 = 31\,500$
  \end{itemize}

  Intersection :
  \[
    x_1 = \max(100,120)=120,\; x_2 = \min(250,270)=250 \;\Rightarrow\; \Delta x = 130
  \]
  \[
    y_1 = \max(150,160)=160,\; y_2 = \min(350,370)=350 \;\Rightarrow\; \Delta y = 190
  \]
  \[
    \text{Aire}_{\text{inter}} = 130 \times 190 = 24\,700
  \]
  Union :
  \[
    \text{Aire}_{\text{union}} = 30\,000 + 31\,500 - 24\,700 = 36\,800
  \]
  \[
    \IoU = \frac{24\,700}{36\,800} \approx \mathbf{0{,}671}
  \]
  La boîte prédite chevauche $67{,}1\,\%$ de la boîte réelle.
  Avec un seuil de $0{,}50$, cette détection est comptée comme \textbf{TP}.
\end{examplebox}

% ─────────────────────────────────────────────────────────────────────────────
\section{GIoU — Generalized IoU}

\begin{definitionbox}{Définition}
  Le \textbf{GIoU} étend l'IoU en pénalisant la distance entre les boîtes
  lorsqu'elles \emph{ne se chevauchent pas}, résolvant le problème des gradients
  nuls de l'IoU standard.
\end{definitionbox}

\begin{formulebox}{Formule}
  \[
    \text{GIoU} = \IoU - \frac{|C \setminus (B_{\text{pred}} \cup B_{\text{gt}})|}{|C|}
  \]
  où $C$ est la \textbf{plus petite boîte englobante convexe} contenant
  $B_{\text{pred}}$ et $B_{\text{gt}}$.
\end{formulebox}

\textbf{Propriétés :}
\[
  -1 \leq \text{GIoU} \leq 1
\]
\begin{itemize}
  \item $\text{GIoU} = 1$ : correspondance parfaite.
  \item $\text{GIoU} \to -1$ : boîtes très éloignées.
  \item $\text{GIoU} = \IoU$ si les boîtes se chevauchent entièrement.
\end{itemize}

\begin{examplebox}{Calcul numérique — boîtes sans chevauchement}
  \begin{itemize}
    \item $B_{\text{pred}} = (0,\,0,\,100,\,100)$ — Aire = $10\,000$
    \item $B_{\text{gt}} = (150,\,150,\,250,\,250)$ — Aire = $10\,000$
    \item $\text{Intersection} = 0$ (aucun chevauchement)
    \item $\IoU = 0$
    \item Boîte $C = (0,\,0,\,250,\,250)$ — Aire$(C) = 62\,500$
    \item Aire$(C) -$ Aire$(\text{union}) = 62\,500 - 20\,000 = 42\,500$
  \end{itemize}
  \[
    \text{GIoU} = 0 - \frac{42\,500}{62\,500} = -\mathbf{0{,}68}
  \]
  Le GIoU de $-0{,}68$ indique que les boîtes sont très éloignées, ce que
  l'IoU = 0 ne pouvait pas quantifier.
\end{examplebox}

% ─────────────────────────────────────────────────────────────────────────────
\section{DIoU — Distance IoU}

\begin{definitionbox}{Définition}
  Le \textbf{DIoU} ajoute une pénalité sur la \textbf{distance entre les centres}
  des deux boîtes, en plus de l'IoU. Il favorise une convergence plus rapide que
  le GIoU.
\end{definitionbox}

\begin{formulebox}{Formule}
  \[
    \text{DIoU} = \IoU - \frac{\rho^2(b,\, b^{gt})}{c^2}
  \]
  \begin{itemize}
    \item $\rho(b, b^{gt})$ : distance euclidienne entre les centres $b$ et
          $b^{gt}$ des deux boîtes.
    \item $c$ : longueur de la diagonale de la plus petite boîte englobante $C$.
  \end{itemize}
\end{formulebox}

\begin{examplebox}{Calcul numérique}
  \begin{itemize}
    \item $B_{\text{pred}} = (50,\,50,\,150,\,150)$
          $\;\Rightarrow\;$ centre $b = (100, 100)$
    \item $B_{\text{gt}} = (80,\,80,\,180,\,180)$
          $\;\Rightarrow\;$ centre $b^{gt} = (130, 130)$
    \item Intersection : $x_1=80, x_2=150, y_1=80, y_2=150$
          $\;\Rightarrow\;$ Aire $= 70 \times 70 = 4\,900$
    \item Aire$_{\text{pred}} = 10\,000$, Aire$_{\text{gt}} = 10\,000$
    \item Aire$_{\text{union}} = 10\,000 + 10\,000 - 4\,900 = 15\,100$
    \item $\IoU = 4\,900 / 15\,100 \approx 0{,}3245$
    \item $\rho^2 = (130-100)^2 + (130-100)^2 = 900 + 900 = 1\,800$
    \item $C = (50,\,50,\,180,\,180)$
          $\;\Rightarrow\;$ $c^2 = 130^2 + 130^2 = 33\,800$
  \end{itemize}
  \[
    \text{DIoU} = 0{,}3245 - \frac{1\,800}{33\,800} \approx 0{,}3245 - 0{,}0533 \approx \mathbf{0{,}271}
  \]
\end{examplebox}

% ─────────────────────────────────────────────────────────────────────────────
\section{CIoU — Complete IoU (Loss par défaut de YOLOv8)}

\begin{definitionbox}{Définition}
  Le \textbf{CIoU} est la métrique de recoupement la plus complète : il intègre
  simultanément l'\textbf{IoU}, la \textbf{distance entre centres} et le
  \textbf{ratio d'aspect} (rapport largeur/hauteur). C'est la fonction de perte
  de localisation native de YOLOv8.
\end{definitionbox}

\begin{formulebox}{Formule}
  \[
    \text{CIoU} = \IoU - \frac{\rho^2(b,\, b^{gt})}{c^2} - \alpha v
  \]
  où :
  \[
    v = \frac{4}{\pi^2} \left(\arctan\frac{w^{gt}}{h^{gt}}
        - \arctan\frac{w}{h}\right)^2
  \]
  \[
    \alpha = \frac{v}{(1 - \IoU) + v}
  \]
  \begin{itemize}
    \item $v$ : mesure la cohérence du ratio d'aspect entre prédiction et GT.
    \item $\alpha$ : coefficient de pondération adaptatif.
    \item $w, h$ : largeur et hauteur de la boîte prédite.
    \item $w^{gt}, h^{gt}$ : largeur et hauteur de la boîte GT.
  \end{itemize}
\end{formulebox}

\begin{examplebox}{Calcul numérique}
  En reprenant l'exemple précédent et ajoutant les dimensions :
  \begin{itemize}
    \item $B_{\text{pred}} : w=100,\; h=100$ (carré)
    \item $B_{\text{gt}} : w=100,\; h=100$ (carré)
    \item $\IoU \approx 0{,}3245$, termes DIoU calculés ci-dessus.
  \end{itemize}
  \[
    v = \frac{4}{\pi^2}\left(\arctan\frac{100}{100} - \arctan\frac{100}{100}\right)^2
      = \frac{4}{\pi^2}(0)^2 = 0
  \]
  Les deux boîtes ayant le \emph{même ratio d'aspect}, $v=0$ et donc
  $\text{CIoU} = \text{DIoU} \approx \mathbf{0{,}271}$.

  \vspace{0.3cm}
  Cas avec ratio d'aspect différent ($B_{\text{pred}} : w=100, h=150$;
  $B_{\text{gt}} : w=100, h=100$) :
  \[
    v = \frac{4}{\pi^2}\left(\arctan(1) - \arctan\tfrac{2}{3}\right)^2
      = \frac{4}{\pi^2}(0{,}7854 - 0{,}5880)^2 \approx \frac{4}{9{,}87}(0{,}0390)
      \approx 0{,}0158
  \]
  \[
    \alpha = \frac{0{,}0158}{(1-0{,}3245)+0{,}0158} = \frac{0{,}0158}{0{,}6913} \approx 0{,}0229
  \]
  \[
    \text{CIoU} \approx 0{,}3245 - 0{,}0533 - 0{,}0229 \times 0{,}0158 \approx \mathbf{0{,}271}
  \]
\end{examplebox}

% ════════════════════════════════════════════════════════════════════════════════
%  CHAPITRE 3 — FONCTIONS DE PERTE
% ════════════════════════════════════════════════════════════════════════════════
\chapter{Fonctions de Perte (Loss Functions)}

Les fonctions de perte guident l'optimisation du modèle pendant l'entraînement.
YOLOv8 dispose de trois pertes distinctes, affichées séparément sur les courbes
d'apprentissage (train et val).

% ─────────────────────────────────────────────────────────────────────────────
\section{Box Loss}

\begin{definitionbox}{Définition}
  La \textbf{Box Loss} quantifie l'erreur de régression des coordonnées des
  boîtes englobantes. YOLOv8 utilise la CIoU Loss (ou DFL + CIoU selon la
  configuration).
\end{definitionbox}

\begin{formulebox}{Formule (CIoU Loss)}
  \[
    \mathcal{L}_{\text{box}} = 1 - \text{CIoU}(B_{\text{pred}},\, B_{\text{gt}})
  \]
\end{formulebox}

\textbf{Interprétation des courbes :}
\begin{itemize}
  \item $\mathcal{L}_{\text{box}} \to 0$ : les boîtes prédites coïncident
        parfaitement avec les boîtes réelles.
  \item Si \texttt{val/box\_loss} $\gg$ \texttt{train/box\_loss} : surapprentissage
        de la localisation.
\end{itemize}

\begin{examplebox}{Exemple sur un batch}
  Soit un batch de 3 détections avec leurs CIoU :
  \begin{center}
  \begin{tabular}{cccc}
    \toprule
    \textbf{Détection} & \textbf{CIoU} & \textbf{Box Loss} & \textbf{Pondération} \\
    \midrule
    Det-1 & 0.85 & $1-0.85=0.15$ & $\times 1.0$ \\
    Det-2 & 0.60 & $1-0.60=0.40$ & $\times 1.0$ \\
    Det-3 & 0.30 & $1-0.30=0.70$ & $\times 1.0$ \\
    \midrule
    \textbf{Moyenne} & — & $\mathbf{0.417}$ & — \\
    \bottomrule
  \end{tabular}
  \end{center}
\end{examplebox}

% ─────────────────────────────────────────────────────────────────────────────
\section{Classification Loss (Cls Loss)}

\begin{definitionbox}{Définition}
  La \textbf{Cls Loss} (ou Classification Loss) mesure l'erreur d'attribution
  des classes aux objets détectés. YOLOv8 utilise la
  \textbf{Binary Cross-Entropy (BCE)} avec sigmoid, appliquée de façon
  indépendante à chaque classe (approche multi-label).
\end{definitionbox}

\begin{formulebox}{Formule}
  Pour chaque classe $c$ et chaque prédiction :
  \[
    \mathcal{L}_{\text{cls}} = -\bigl[y_c \log(\hat{p}_c) +
    (1 - y_c)\log(1 - \hat{p}_c)\bigr]
  \]
  \begin{itemize}
    \item $y_c \in \{0, 1\}$ : label réel (1 si l'objet appartient à la classe $c$)
    \item $\hat{p}_c \in [0,1]$ : probabilité prédite par le modèle pour la classe $c$
  \end{itemize}
\end{formulebox}

\begin{examplebox}{Calcul numérique — 3 classes}
  Objet de classe ``voiture'' ($y_1=1$) avec probabilités prédites :
  \begin{center}
  \begin{tabular}{lcccc}
    \toprule
    \textbf{Classe} & $y_c$ & $\hat{p}_c$ & $\mathcal{L}_c$ \\
    \midrule
    Voiture  & 1 & 0.85 & $-(1 \times \log 0.85) = 0.163$ \\
    Piéton   & 0 & 0.10 & $-(1 \times \log 0.90) = 0.105$ \\
    Vélo     & 0 & 0.05 & $-(1 \times \log 0.95) = 0.051$ \\
    \midrule
    \textbf{Total} & — & — & $\mathbf{0.319}$ \\
    \bottomrule
  \end{tabular}
  \end{center}
  La loss est faible car le modèle attribue correctement une haute probabilité à
  ``voiture'' et de faibles probabilités aux autres classes.
\end{examplebox}

% ─────────────────────────────────────────────────────────────────────────────
\section{DFL Loss — Distribution Focal Loss}

\begin{definitionbox}{Définition}
  La \textbf{DFL Loss} est une innovation de YOLOv8 : au lieu de prédire
  directement les coordonnées scalaires des bords des boîtes, le modèle prédit
  une \textbf{distribution de probabilité} sur un ensemble discret de valeurs
  possibles. La DFL Loss optimise cette distribution.
\end{definitionbox}

\begin{formulebox}{Formule}
  \[
    \text{DFL}(S_i,\, S_{i+1}) =
    -\Bigl[\,(y_{i+1} - y)\,\log(S_i) + (y - y_i)\,\log(S_{i+1})\,\Bigr]
  \]
  \begin{itemize}
    \item $y$ : valeur cible du bord (continue)
    \item $y_i, y_{i+1}$ : deux valeurs discrètes encadrant $y$
          (i.e.\ $y_i \leq y < y_{i+1}$)
    \item $S_i, S_{i+1}$ : probabilités prédites pour $y_i$ et $y_{i+1}$
          (issues du softmax)
  \end{itemize}
\end{formulebox}

\textbf{Intuition :} le modèle apprend que le bord gauche d'une boîte est à
``environ 45 pixels'', avec une distribution gaussienne centrée sur 45 — plus
robuste qu'une valeur scalaire précise.

\begin{examplebox}{Calcul numérique}
  Supposons $y = 45{,}3$ (position cible d'un bord), $y_i = 45$, $y_{i+1} = 46$,
  et le modèle prédit $S_{45} = 0{,}72$ et $S_{46} = 0{,}28$ :
  \[
    \text{DFL} = -\bigl[(46-45{,}3)\log(0{,}72) + (45{,}3-45)\log(0{,}28)\bigr]
  \]
  \[
    = -\bigl[0{,}7 \times (-0{,}3285) + 0{,}3 \times (-1{,}2730)\bigr]
  \]
  \[
    = -\bigl[-0{,}2300 - 0{,}3819\bigr] = \mathbf{0{,}612}
  \]
  Si le modèle était parfait ($S_{45}=0{,}7$, $S_{46}=0{,}3$) :
  \[
    \text{DFL} = -\bigl[0{,}7 \times \log(0{,}7) + 0{,}3 \times \log(0{,}3)\bigr]
               \approx -\bigl[-0{,}2490 - 0{,}3612\bigr] = 0{,}610
  \]
\end{examplebox}

\begin{notebox}
  La DFL Loss est combinée à la CIoU Loss dans la Box Loss totale de YOLOv8 :
  $\mathcal{L}_{\text{box}} = \lambda_{\text{CIoU}} \cdot \mathcal{L}_{\text{CIoU}}
  + \lambda_{\text{DFL}} \cdot \mathcal{L}_{\text{DFL}}$
\end{notebox}

% ════════════════════════════════════════════════════════════════════════════════
%  CHAPITRE 4 — MÉTRIQUES OPÉRATIONNELLES
% ════════════════════════════════════════════════════════════════════════════════
\chapter{Métriques Opérationnelles}

\section{FPS — Frames Per Second}

\begin{definitionbox}{Définition}
  Les \textbf{FPS} mesurent le nombre d'images traitées par seconde. Indicateur
  clé pour les applications temps réel.
\end{definitionbox}

\begin{formulebox}{Formule}
  \[
    \text{FPS} = \frac{1}{\text{Latence (s)}} = \frac{\text{Nombre d'images}}{\text{Temps total (s)}}
  \]
\end{formulebox}

\begin{examplebox}{Comparaison des modèles YOLOv8}
  \begin{center}
  \begin{tabular}{lcccc}
    \toprule
    \textbf{Modèle} & \textbf{Latence (ms)} & \textbf{FPS} & \textbf{mAP50} & \textbf{Params} \\
    \midrule
    YOLOv8n & 1.47 & 680 & 37.3 & 3.2M \\
    YOLOv8s & 2.66 & 376 & 44.9 & 11.2M \\
    YOLOv8m & 5.86 & 171 & 50.2 & 25.9M \\
    YOLOv8l & 9.06 & 110 & 52.9 & 43.7M \\
    YOLOv8x & 14.37 & 70  & 53.9 & 68.2M \\
    \bottomrule
  \end{tabular}
  \end{center}
  \textit{Source : Ultralytics benchmark, GPU Tesla T4, image 640$\times$640.}
\end{examplebox}

\section{GFLOPs — Giga Floating Point Operations}

\begin{definitionbox}{Définition}
  Les \textbf{GFLOPs} mesurent la complexité computationnelle d'une passe avant
  (forward pass) du modèle. Indicateur de la charge de calcul requise.
\end{definitionbox}

\begin{formulebox}{Formule approximative (couche convolutive)}
  \[
    \text{FLOPs}_{\text{conv}} =
      2 \times C_{\text{in}} \times C_{\text{out}} \times K_w \times K_h \times H_{\text{out}} \times W_{\text{out}}
  \]
  Le facteur 2 représente une multiplication et une addition.
\end{formulebox}

\begin{examplebox}{Exemple de calcul sur une couche}
  Couche convolutive : $C_{\text{in}}=64$, $C_{\text{out}}=128$, $K=3\times3$,
  sortie $H_{\text{out}} \times W_{\text{out}} = 80\times80$ :
  \[
    \text{FLOPs} = 2 \times 64 \times 128 \times 3 \times 3 \times 80 \times 80
                = 2 \times 64 \times 128 \times 9 \times 6400
                \approx 943\,\text{M FLOPs} = 0{,}94\;\text{GFLOPs}
  \]
\end{examplebox}

\section{Paramètres du modèle (Params)}

\begin{definitionbox}{Définition}
  Le nombre de \textbf{paramètres} (poids) détermine la capacité d'apprentissage
  du modèle et son empreinte mémoire.
\end{definitionbox}

\begin{formulebox}{Formule (couche convolutive)}
  \[
    \text{Params}_{\text{conv}} = C_{\text{out}} \times (C_{\text{in}} \times K_w \times K_h + 1)
  \]
  Le $+1$ représente le biais.
\end{formulebox}

\begin{examplebox}{Calcul}
  Pour la couche de l'exemple précédent :
  \[
    \text{Params} = 128 \times (64 \times 3 \times 3 + 1) = 128 \times 577 = 73\,856
    \approx 74\text{k paramètres}
  \]
\end{examplebox}

% ════════════════════════════════════════════════════════════════════════════════
%  CHAPITRE 5 — MÉTRIQUES DIAGNOSTIQUES
% ════════════════════════════════════════════════════════════════════════════════
\chapter{Métriques de Diagnostic}

\section{Confidence Score}

\begin{definitionbox}{Définition}
  Le \textbf{score de confiance} est la probabilité jointe qu'une région contient
  un objet \emph{et} que cet objet appartient à une classe donnée.
\end{definitionbox}

\begin{formulebox}{Formule}
  \[
    \text{Conf} = P(\text{object}) \times P(\text{class} \mid \text{object})
                = P(\text{object}) \times P_c
  \]
\end{formulebox}

\begin{examplebox}{Exemple}
  \begin{itemize}
    \item $P(\text{object}) = 0{,}92$ — le modèle est sûr à 92\% qu'il y a un objet.
    \item $P(\text{voiture} \mid \text{objet}) = 0{,}88$ — parmi les objets, 88\%
          de chances que ce soit une voiture.
    \item $\text{Conf} = 0{,}92 \times 0{,}88 = \mathbf{0{,}81}$
  \end{itemize}
  Cette détection passe le seuil de confiance par défaut de 0{,}25.
\end{examplebox}

\section{Courbe PR (Precision-Recall)}

\begin{definitionbox}{Définition}
  La \textbf{courbe PR} trace la Precision en fonction du Recall lorsqu'on fait
  varier le seuil de confiance de 1{,}0 (très restrictif) à 0{,}0 (très permissif).
\end{definitionbox}

\textbf{Interprétation :}
\begin{itemize}
  \item Courbe proche du coin supérieur droit $(R=1, P=1)$ : modèle excellent.
  \item Aire sous la courbe = \textbf{AP} (voir section précédente).
  \item Chute rapide de la Precision quand le Recall augmente : le modèle produit
        beaucoup de fausses alarmes pour ``capturer'' les derniers objets.
\end{itemize}

\begin{center}
\begin{tikzpicture}
  \begin{axis}[
    width=10cm, height=7cm,
    xlabel={Recall}, ylabel={Precision},
    xmin=0, xmax=1, ymin=0, ymax=1,
    grid=both, grid style={gray!30},
    title={\textbf{Courbe PR — Exemple YOLOv8}},
    legend pos=south west,
    tick label style={font=\small}
  ]
  % Courbe d'un bon modèle
  \addplot[color=yoloblue, thick, smooth] coordinates {
    (0.0,1.00)(0.10,0.98)(0.20,0.96)(0.30,0.93)
    (0.40,0.89)(0.50,0.85)(0.60,0.80)(0.70,0.73)
    (0.80,0.63)(0.90,0.50)(1.00,0.30)
  };
  \addlegendentry{YOLOv8m (AP=0.81)}
  % Courbe d'un modèle moyen
  \addplot[color=yoloorange, thick, smooth, dashed] coordinates {
    (0.0,0.90)(0.10,0.84)(0.20,0.78)(0.30,0.71)
    (0.40,0.64)(0.50,0.57)(0.60,0.50)(0.70,0.42)
    (0.80,0.33)(0.90,0.22)(1.00,0.10)
  };
  \addlegendentry{YOLOv8n (AP=0.56)}
  % Ligne de référence (modèle aléatoire)
  \addplot[color=gray, thin, dotted] coordinates {(0,0.5)(1,0.5)};
  \addlegendentry{Aléatoire}
  \end{axis}
\end{tikzpicture}
\end{center}

\section{Courbe F1-Confidence}

\begin{definitionbox}{Définition}
  La \textbf{courbe F1-Confidence} affiche le F1-Score en fonction du seuil de
  confiance. Elle permet d'identifier le seuil optimal à utiliser lors de
  l'inférence.
\end{definitionbox}

\begin{examplebox}{Tableau d'analyse}
  \begin{center}
  \begin{tabular}{ccccc}
    \toprule
    \textbf{Seuil conf.} & \textbf{P} & \textbf{R} & \textbf{F1} & \textbf{Recommandation} \\
    \midrule
    0.10 & 0.61 & 0.95 & 0.745 & Trop permissif \\
    0.25 & 0.72 & 0.88 & 0.792 & Défaut YOLOv8 \\
    0.40 & 0.81 & 0.79 & \textbf{0.800} & \textbf{Optimal} \\
    0.55 & 0.89 & 0.68 & 0.771 & Précis, moins de rappel \\
    0.70 & 0.94 & 0.52 & 0.671 & Très restrictif \\
    0.85 & 0.98 & 0.31 & 0.471 & Trop restrictif \\
    \bottomrule
  \end{tabular}
  \end{center}
  Le seuil optimal est 0{,}40, maximisant le F1-Score à 0{,}800.
\end{examplebox}

\section{Matrice de Confusion}

\begin{definitionbox}{Définition}
  La \textbf{matrice de confusion} est un tableau $N \times N$ (où $N$ est le
  nombre de classes + 1 pour ``background'') croisant classes prédites et classes
  réelles. Elle révèle les confusions inter-classes et les taux de détection
  manquée.
\end{definitionbox}

\begin{examplebox}{Matrice de confusion — 3 classes}
  Normalisée par rapport au total des GT par classe :
  \begin{center}
  \begin{tabular}{l|ccc|c}
    \toprule
    & \multicolumn{3}{c|}{\textbf{Prédit}} & \\
    \textbf{Réel} & Voiture & Piéton & Vélo & Manqué \\
    \midrule
    Voiture & \cellcolor{yologreen!40}\textbf{0.89} & 0.03 & 0.02 & 0.06 \\
    Piéton  & 0.05 & \cellcolor{yologreen!40}\textbf{0.76} & 0.08 & 0.11 \\
    Vélo    & 0.02 & 0.12 & \cellcolor{yologreen!40}\textbf{0.71} & 0.15 \\
    \bottomrule
  \end{tabular}
  \end{center}

  \textit{Lecture :}
  \begin{itemize}
    \item 89\% des voitures réelles sont détectées comme ``voiture''.
    \item 12\% des vélos sont confondus avec des piétons.
    \item 15\% des vélos sont manqués (FN).
  \end{itemize}
\end{examplebox}

% ════════════════════════════════════════════════════════════════════════════════
%  CHAPITRE 6 — TABLEAU RÉCAPITULATIF
% ════════════════════════════════════════════════════════════════════════════════
\chapter{Tableau Récapitulatif Exhaustif}

\begin{longtable}{>{\bfseries}p{3.2cm} p{3.8cm} p{3.5cm} p{3.5cm}}
  \toprule
  \textbf{Métrique} & \textbf{Formule} & \textbf{Plage} & \textbf{Rôle} \\
  \midrule
  \endfirsthead
  \toprule
  \textbf{Métrique} & \textbf{Formule} & \textbf{Plage} & \textbf{Rôle} \\
  \midrule
  \endhead
  \bottomrule
  \endlastfoot

  Precision (P)
    & $\frac{TP}{TP+FP}$
    & $[0, 1]$, $\uparrow$
    & Taux de détections correctes \\[4pt]

  Recall (R)
    & $\frac{TP}{TP+FN}$
    & $[0, 1]$, $\uparrow$
    & Taux d'objets réels détectés \\[4pt]

  F1-Score
    & $2 \cdot \frac{P \cdot R}{P+R}$
    & $[0, 1]$, $\uparrow$
    & Équilibre P/R, choix du seuil \\[4pt]

  AP (par classe)
    & $\int_0^1 P(R)\,dR$
    & $[0, 1]$, $\uparrow$
    & Précision globale d'une classe \\[4pt]

  mAP50
    & $\frac{1}{N}\sum AP_i^{@0.5}$
    & $[0, 1]$, $\uparrow$
    & Métrique principale YOLOv8 \\[4pt]

  mAP50-95
    & $\frac{1}{10}\sum_{t} \mAP_t$
    & $[0, 1]$, $\uparrow$
    & Standard COCO, localisation fine \\[4pt]

  \midrule

  IoU
    & $\frac{|B_p \cap B_g|}{|B_p \cup B_g|}$
    & $[0, 1]$, $\uparrow$
    & Qualité de localisation \\[4pt]

  GIoU
    & $\IoU - \frac{|C \setminus \text{union}|}{|C|}$
    & $[-1, 1]$, $\uparrow$
    & IoU + pénalité distance \\[4pt]

  DIoU
    & $\IoU - \frac{\rho^2}{c^2}$
    & $[-1, 1]$, $\uparrow$
    & IoU + distance des centres \\[4pt]

  CIoU
    & $\text{DIoU} - \alpha v$
    & $[-1, 1]$, $\uparrow$
    & Loss de localisation YOLOv8 \\[4pt]

  \midrule

  Box Loss
    & $1 - \text{CIoU}$
    & $[0, +\infty)$, $\downarrow$
    & Erreur de régression des boîtes \\[4pt]

  Cls Loss
    & BCE par classe
    & $[0, +\infty)$, $\downarrow$
    & Erreur de classification \\[4pt]

  DFL Loss
    & $-[(\cdot)\log S_i + (\cdot)\log S_{i+1}]$
    & $[0, +\infty)$, $\downarrow$
    & Précision des bords (distribution) \\[4pt]

  \midrule

  FPS
    & $1/\text{latence}$
    & $\mathbb{R}^+$, $\uparrow$
    & Vitesse temps réel \\[4pt]

  GFLOPs
    & Opérations flottantes
    & $\mathbb{R}^+$, $\downarrow$
    & Charge computationnelle \\[4pt]

  Params (M)
    & Nb de poids du modèle
    & $\mathbb{R}^+$, contexte
    & Taille et capacité du modèle \\[4pt]

  Confidence
    & $P(\text{obj}) \times P_c$
    & $[0, 1]$, $\uparrow$
    & Certitude d'une détection \\[4pt]

  Courbe PR
    & $P = f(R)$
    & Visualisation
    & Trade-off P/R global \\[4pt]

  Courbe F1-Conf
    & $F1 = f(\text{conf})$
    & Visualisation
    & Seuil de confiance optimal \\[4pt]

  Matrice confusion
    & $N \times N$ classes
    & Visualisation
    & Confusions inter-classes \\

\end{longtable}

\textit{$\uparrow$ : une valeur plus élevée est meilleure. $\downarrow$ : une
valeur plus basse est meilleure.}

% ════════════════════════════════════════════════════════════════════════════════
%  CHAPITRE 7 — CODE PYTHON DE RÉFÉRENCE
% ════════════════════════════════════════════════════════════════════════════════
\chapter{Implémentation Python de Référence}

\begin{lstlisting}[style=python, caption={Calcul des métriques de base avec NumPy},
  label={lst:metrics_numpy}]
import numpy as np

# ──────────────────────────────────────────────────────────────────────────────
# 1. Métriques de classification
# ──────────────────────────────────────────────────────────────────────────────
TP, FP, FN = 6, 3, 2

precision = TP / (TP + FP)
recall    = TP / (TP + FN)
f1        = 2 * precision * recall / (precision + recall)

print(f"Precision : {precision:.4f}")  # 0.6667
print(f"Recall    : {recall:.4f}")     # 0.7500
print(f"F1-Score  : {f1:.4f}")         # 0.7059

# ──────────────────────────────────────────────────────────────────────────────
# 2. Calcul de l'IoU
# ──────────────────────────────────────────────────────────────────────────────
def compute_iou(box_pred, box_gt):
    """
    box format: (x1, y1, x2, y2)
    """
    x1 = max(box_pred[0], box_gt[0])
    y1 = max(box_pred[1], box_gt[1])
    x2 = min(box_pred[2], box_gt[2])
    y2 = min(box_pred[3], box_gt[3])

    inter = max(0, x2 - x1) * max(0, y2 - y1)
    area_pred = (box_pred[2]-box_pred[0]) * (box_pred[3]-box_pred[1])
    area_gt   = (box_gt[2]-box_gt[0])   * (box_gt[3]-box_gt[1])
    union = area_pred + area_gt - inter

    return inter / union if union > 0 else 0.0

iou = compute_iou((100,150,250,350), (120,160,270,370))
print(f"IoU       : {iou:.4f}")  # 0.6712

# ──────────────────────────────────────────────────────────────────────────────
# 3. Calcul de l'AP (méthode des 11 points)
# ──────────────────────────────────────────────────────────────────────────────
def compute_ap_11pts(precisions, recalls):
    """
    precisions, recalls : listes triées par confiance décroissante
    """
    ap = 0.0
    for t in np.linspace(0, 1, 11):
        # Precision max pour Recall >= t
        p_interp = max(
            [p for p, r in zip(precisions, recalls) if r >= t],
            default=0
        )
        ap += p_interp / 11
    return ap

prec = [1.0, 1.0, 0.667, 0.75, 0.6, 0.667, 0.571, 0.625, 0.556, 0.5]
rec  = [0.2, 0.4, 0.4,   0.6,  0.6, 0.8,   0.8,   1.0,   1.0,   1.0]
ap   = compute_ap_11pts(prec, rec)
print(f"AP (11pts): {ap:.4f}")  # ≈ 0.5848

# ──────────────────────────────────────────────────────────────────────────────
# 4. Entraînement YOLOv8 avec Ultralytics
# ──────────────────────────────────────────────────────────────────────────────
from ultralytics import YOLO

model = YOLO("yolov8m.pt")           # Charger le modèle pré-entraîné

results = model.train(
    data   = "coco128.yaml",         # Fichier de configuration du dataset
    epochs = 100,
    imgsz  = 640,
    batch  = 16,
    device = 0                       # GPU 0
)

# 5. Évaluation
metrics = model.val(data="coco128.yaml", imgsz=640)
print(f"mAP50    : {metrics.box.map50:.4f}")
print(f"mAP50-95 : {metrics.box.map:.4f}")
print(f"Precision: {metrics.box.mp:.4f}")
print(f"Recall   : {metrics.box.mr:.4f}")
\end{lstlisting}

% ════════════════════════════════════════════════════════════════════════════════
%  CONCLUSION
% ════════════════════════════════════════════════════════════════════════════════
\chapter{Conclusion et Guide Pratique}

\section{Quelle métrique utiliser selon l'objectif ?}

\begin{center}
\begin{tabular}{p{4.5cm} p{5cm} p{4cm}}
  \toprule
  \textbf{Objectif} & \textbf{Métrique prioritaire} & \textbf{Secondaire} \\
  \midrule
  Comparer des modèles & mAP50-95 (COCO) & mAP50 \\
  Déploiement rapide & FPS / Latence & GFLOPs \\
  Minimiser fausses alarmes & Precision & F1-Score \\
  Minimiser objets manqués & Recall & F1-Score \\
  Équilibre général & F1-Score / mAP50 & Courbe PR \\
  Analyser les erreurs & Matrice de confusion & Courbe F1-Conf \\
  Localisation précise & mAP50-95 / CIoU & DIoU \\
  Embarqué / Edge AI & Params (M) / Size & GFLOPs \\
  \bottomrule
\end{tabular}
\end{center}

\section{Seuils de référence YOLOv8 (COCO)}

\begin{center}
\begin{tabular}{lcccc}
  \toprule
  \textbf{Modèle} & \textbf{mAP50} & \textbf{mAP50-95} & \textbf{Params} & \textbf{FPS (T4)} \\
  \midrule
  YOLOv8n & 52.9 & 37.3 & 3.2M  & 680 \\
  YOLOv8s & 61.8 & 44.9 & 11.2M & 376 \\
  YOLOv8m & 67.2 & 50.2 & 25.9M & 171 \\
  YOLOv8l & 69.8 & 52.9 & 43.7M & 110 \\
  YOLOv8x & 71.0 & 53.9 & 68.2M & 70  \\
  \bottomrule
\end{tabular}
\end{center}

\section{Checklist d'évaluation YOLOv8}

\begin{tcolorbox}[colback=lightblue, colframe=yoloblue,
  title={\textbf{Checklist : évaluer un modèle YOLOv8}}, breakable]
  \begin{itemize}
    \item[$\square$] Vérifier la convergence des \textbf{3 losses} (box, cls, dfl)
          sur train \emph{et} val.
    \item[$\square$] Contrôler l'absence de \textbf{surapprentissage} :
          val\_loss $\leq 1{,}5 \times$ train\_loss.
    \item[$\square$] Comparer \textbf{mAP50} et \textbf{mAP50-95} : écart $> 0{,}25$
          = localisation imprécise.
    \item[$\square$] Analyser la \textbf{matrice de confusion} pour détecter les
          confusions inter-classes.
    \item[$\square$] Utiliser la \textbf{courbe F1-Confidence} pour choisir le
          seuil de confiance optimal.
    \item[$\square$] Vérifier \textbf{AP par classe} : classes avec AP $< 0{,}5$
          nécessitent plus de données.
    \item[$\square$] Mesurer les \textbf{FPS} sur le matériel cible.
  \end{itemize}
\end{tcolorbox}

\bigskip
\begin{notebox}
  Ce document couvre l'ensemble des métriques utilisées par YOLOv8 dans le
  contexte de la détection et de la reconnaissance d'objets. Pour des cas
  d'usage spécifiques (segmentation, pose estimation), des métriques
  supplémentaires comme le \textbf{Mask AP} ou l'**OKS (Object Keypoint
  Similarity)** entrent en jeu, mais restent des extensions des principes
  présentés ici.
\end{notebox}

\end{document}
